\chapter*{ВВЕДЕНИЕ}
\addcontentsline{toc}{chapter}{ВВЕДЕНИЕ}

\textbf{Актуальность темы.}
Робототехника является одной из наиболее динамично развивающихся областей современной техники.
Проектирование, верификация и тестирование робототехнических систем требуют значительных временных и финансовых затрат при работе с физическим оборудованием.
Программные симуляторы позволяют перенести большую часть этих работ в виртуальную среду, существенно сокращая цикл разработки и снижая риски.

Физическое моделирование — ключевой компонент любой симуляционной платформы.
Его задача состоит в воспроизведении динамики твёрдых тел, контактных взаимодействий и кинематических ограничений с точностью, достаточной для верификации алгоритмов управления.
Требования к производительности при этом постоянно растут: задачи обучения с подкреплением предполагают запуск тысяч параллельных симуляций, а построение цифровых двойников требует работы в режиме реального времени.
Графические процессоры (GPU), обладающие тысячами вычислительных ядер, открывают принципиально новые возможности для ускорения расчётов физической динамики.

В условиях ограничений на использование зарубежных программных продуктов актуальность создания отечественных инструментов робототехнической симуляции существенно возрастает.
Разработка программного комплекса «Модейлер» направлена на создание полноценной платформы для моделирования и тестирования робототехнических систем, ориентированной на задачи импортонезависимой инженерной разработки.
Физический движок с GPU-ускорением является одной из центральных подсистем этого комплекса.

\textbf{Степень научной разработанности.}
Задача физического моделирования твёрдых тел в робототехнике изучается начиная с 1990-х годов.
Основополагающие работы Д.~Барафа и А.~Уиткина~\cite{baraff1997physically} заложили теоретическую базу для моделирования контактной динамики.
На сегодняшний день существует ряд широко используемых библиотек физического моделирования: Bullet Physics~\cite{coumans2015bullet}, NVIDIA PhysX~\cite{physx2023}, MuJoCo~\cite{todorov2012mujoco}, ODE~\cite{smith2005ode} и Rapier~\cite{crozet2021rapier}.
В 2024 году компания NVIDIA анонсировала физический движок Newton~\cite{nvidia2024newton}, ориентированный на GPU-ускоренную симуляцию для задач робототехники.

Тем не менее перечисленные решения, как правило, не объединяют в одном продукте удобную среду визуального моделирования, высокоточную физическую симуляцию с масштабным GPU-ускорением и глубокую интеграцию с инструментами разработки робототехнических систем.
Кроме того, большинство из них распространяются на условиях лицензий, ограничивающих коммерческое или государственное применение.

\textbf{Цель исследования} — разработка подсистемы физического моделирования с поддержкой GPU-ускорения для применения в программном комплексе робототехнической симуляции «Модейлер».

\textbf{Задачи исследования:}
\begin{enumerate}
  \item провести анализ задач физической симуляции в робототехнике и требований к физическим движкам;
  \item выполнить сравнительный обзор существующих систем физического моделирования;
  \item определить функциональные и нефункциональные требования к разрабатываемой подсистеме;
  \item обосновать выбор алгоритмов динамики твёрдых тел, обнаружения столкновений и решения ограничений;
  \item разработать архитектуру физического движка и модуля GPU-ускорения;
  \item реализовать ядро физического движка и GPU-вычислительный модуль на основе Vulkan Compute;
  \item провести верификацию физической корректности и оценить производительность разработанной подсистемы.
\end{enumerate}

\textbf{Объект исследования} — программный комплекс моделирования робототехнических систем «Модейлер».

\textbf{Предмет исследования} — алгоритмы и программные средства физического моделирования динамики твёрдых тел с применением вычислений на графических процессорах.

\textbf{Практическая значимость.}
Разработанная подсистема является компонентом программного комплекса «Модейлер» и может применяться для:
верификации алгоритмов управления роботами в режиме software-in-the-loop;
построения цифровых двойников промышленных и мобильных робототехнических установок;
обучения нейросетевых систем управления методами обучения с подкреплением.
Применение GPU-ускорения позволяет многократно увеличить число одновременно моделируемых систем по сравнению с CPU-реализацией.

\textbf{Структура работы.}
Диссертация состоит из введения, пяти глав, заключения, списка использованных источников и приложений.

В первой главе проведён анализ предметной области: рассмотрены задачи физической симуляции в робототехнике, выполнен обзор существующих решений и сформулированы требования к разрабатываемой системе.

Вторая глава посвящена теоретическим основам: изложены уравнения механики твёрдого тела, рассмотрены алгоритмы обнаружения столкновений, методы численного интегрирования и модели контактной динамики.

В третьей главе представлены архитектура программного комплекса «Модейлер» и детальное проектирование физического движка и GPU-вычислительного модуля.

Четвёртая глава описывает реализацию: ядро физического движка, GPU-ускорение на основе Vulkan Compute и программные интерфейсы.

Пятая глава содержит результаты тестирования и оценку производительности.
